โรคหัวใจและหลอดเลือด (CVDs) เป็นสาเหตุสำคัญของการเสียชีวิตทั่วโลก โครงงานปริญญานิพนธ์วิศวกรรมศาสตรบัณฑิตนี้ นำเสนอการออกแบบและพัฒนาระบบตรวจคัดกรองความเสี่ยงสุขภาพหัวใจและหลอดเลือดแบบไม่รุกรานโดยไม่พึ่งพาฮาร์ดแวร์ภายนอก พร้อมการจำลองพลวัตหัวใจดิจิทัลทวิน (In-Silico Digital Twin) ดำเนินงานโดยทีมนักศึกษา 2 คน ระบบทำการบันทึกสัญญาณแรงสั่นสะเทือนเชิงกลของหัวใจที่กระดูกอก (Sternal Seismocardiography: SCG) ตามแนวแกน Dorso-Ventral ด้วยเซนเซอร์วัดความเร่ง 3 แกน (MEMS Accelerometer) ที่ติดตั้งอยู่ภายในสมาร์ทโฟนทั่วไปขณะผู้ใช้อยู่ในท่านอนหงายราบ

การประมวลผลสัญญาณดิจิทัลประยุกต์ใช้วงจรกรองแถบความถี่ Butterworth อันดับ 4 แบบ Zero-Phase (0.5--40.0 Hz) ร่วมกับเกณฑ์ตรวจสอบพลังงานทางสรีรวิทยา ($0.005\text{ m/s}^2 \le \sigma_{raw} \le 0.50\text{ m/s}^2$) และระบบแจ้งเตือนให้ผู้ใช้อยู่นิ่งแบบเรียลไทม์ การสกัดตำแหน่งยอดคลื่นการบีบตัวของหัวใจ (Systolic Peak) ใช้เทคนิค Moving-RMS Energy Envelope และ Adaptive Threshold เพื่อสร้างกราฟช่วงเวลาการเต้นของหัวใจ (PPI Tachogram) และสกัดคุณลักษณะความแปรปรวนของอัตราชีพจร (Pulse Rate Variability: PRV) ทั้งในโดเมนเวลาและโดเมนความถี่ เพื่อป้อนเข้าสู่โมเดลการเรียนรู้ของเครื่อง (Machine Learning) ในการจำแนกความผิดปกติของจังหวะการเต้นหัวใจและความเครียดทางระบบประสาทอัตโนมัติ

ผลการทดสอบกับชุดข้อมูลมาตรฐานทางการแพทย์ (PhysioNet MIT-BIH และ WESAD) พบว่าโมเดล XGBoost ให้ความถูกต้อง (Accuracy) ในการจำแนกความผิดปกติของจังหวะหัวใจสูงถึง 89.4\% ความไว (Sensitivity) 88.9\% ความจำเพาะ (Specificity) 89.8\% และ F1-score 0.887 โดยมีความคลาดเคลื่อนสัมบูรณ์เฉลี่ยในการประมาณอัตราชีพจรเพียง 2.14 ครั้งต่อนาที (MAE = 2.14 bpm) เมื่อเทียบกับการตรวจคลื่นไฟฟ้าหัวใจ (ECG) มาตรฐาน และมีความหน่วงเวลาประมวลผลบนอุปกรณ์เพียง 56.1 มิลลิวินาที นอกจากนี้ คุณลักษณะ PRV ยังสามารถนำไปขับเคลื่อนแบบจำลอง Van der Pol Nonlinear Oscillator เพื่อแสดงภาพพฤติกรรมการเสื่อมสภาพของเสถียรภาพจังหวะหัวใจได้อย่างมีประสิทธิภาพ
